\pagebreak

\section{External ZBT RAM interface - zbt\_ram.3pl}

    \index{library!zbt\_ram.3pl}
    This library provides read and write interfaces for zero bus turnaround
    (ZBT) synchronous static RAM.
    
    This library must be incorporated by using a {\bf module} statement -
    \begin{lstlisting}[gobble=4, language=threepl]
       module "zbt_ram.3pl"
    \end{lstlisting}
    
    
    The RAM address and data dimensions are inferred from the dimensions
    of the address and data FPGA pin arguments supplied to the class
    call.

    Multiple read and write ports are allowed. Addresses and data are
    transferred using queue mode variables and simultaneous requests are handled
    by arbitration based on fair access or priority - see below.
    
    The interfaces to a single external RAM chip are incorporated in
    an instance of a class. Reads and writes are provided by modules in
    that class.


    

    \subsection{library modules}
    
        none.

    \subsection{library procedures}
    
        None.

    \subsection{library functions}
    
        None.


    \subsection{library classes}
    
        \subsubsection{class zbt\_ram\label{sec:lczbtram}}
        
            {\bf Prototype:} \\
            \vsp
            \begin{lstlisting}[gobble=12, language=threepl]
                global class
                zbt_ram (
                    var(addr_pins, "[]str"),
                    var(data_pins, "[]str"),
                    str(_we_pin),
                    str(_wea_pin),
                    str(_web_pin),
                    str(_wec_pin),
                    str(_wed_pin),
                    str(_cke_pin),
                    clock(clk),
                    str(slew, "fast"),
                    int(drive, 12),     // 12mA
                    str(tnm, "ZBTPAD")
                ) {}
            \end{lstlisting}

            {\bf Output parameters:} \\
            
            none.

            {\bf Input parameters:} \\
            {\bf addr\_pins} are the RAM address FPGA output pins \\
            {\bf data\_pins} are the RAM data FPGA bi-directional pins \\
            {\bf \_we\_pin} is the RAM write enable FPGA output pin \\
            {\bf \_wea\_pin} is the RAM byte A write enable FPGA output pin \\
            {\bf \_web\_pin} is the RAM byte B write enable FPGA output pin \\
            {\bf \_wec\_pin} is the RAM byte C write enable FPGA output pin \\
            {\bf \_wed\_pin} is the RAM byte D write enable FPGA output pin \\
            {\bf \_cke\_pin} is the RAM clock enable FPGA output pin \\
            {\bf clk} is the clock domain of the RAM interface \\
            {\bf slew} is the output pad slew rate ("fast", "slow") \\
            {\bf drive} is the output pad current drive (2 to 24 mA)\\
            {\bf tnm} is the timing group name for the pads \\
            
            It is presumed that the external RAM chip has a clock source
            which is also input to the FPGA and that {\bf clk} is the
            FPGA version of that clock. If that argument is missing the
            current clock domain is used for the interface logic.
            
            The output pad slew rate and drive current are optional arguments
            which if omitted have the default values shown.
            
            The timing group name {\bf tnm} is applied to all input and output pads
            connected to this interface.
            
 
            The following is an example of creation of a class instance - 
 
            \begin{lstlisting}[gobble=12, language=threepl]
                class(zbt);
                zbt =   zbt_ram(
                            addr_pins=ext_ram_a_pins[0],
                            data_pins=ext_ram_d_pins[0],
                            _we_pin=_ext_ram_we_pins[0],
                            _wea_pin=_ext_ram_wea_pins[0],
                            _web_pin=_ext_ram_web_pins[0],
                            _wec_pin=_ext_ram_wec_pins[0],
                            _wed_pin=_ext_ram_wed_pins[0],
                            _cke_pin=_ext_ram_cke_pins[0],
                            clk=c166
                        );
            \end{lstlisting}

 
            


        \subsubsection{class module read - read from external ZBT memory\\
                       class module readp - read from external ZBT memory with priority}\label{sec:lczbtramread}

            {\bf Library: zbt\_ram.3pl}

            {\bf Prototypes:} \\
            \vsp
            \begin{lstlisting}[gobble=12, language=threepl]
               module
               read (queue(data) <- queue(address))) {}
            \end{lstlisting}
            \begin{lstlisting}[gobble=12, language=threepl]
               module
               readp (queue(data) <- int(pri), queue(address)) {}
            \end{lstlisting}

            {\bf Output parameters:} \\
            {\bf data} is the queue receiving the data read from the memory.

            {\bf Input parameters:} \\
            {\bf pri} is the priority of the read request. \\
            {\bf addr} is the address to be read.

            {\bf description:} \\
            When the output and input queues are available an external memory
            read will be executed, subject to access arbitration. The address
            input queue and the data output queue may be any type, but a
            crude pad/truncate cast to the address and data primitive sizes
            will be performed.
            
            The internal code of this module executes in clock domain {\bf clk}
            (but the call to the module returns with the current clock domain
            restored). If the clock domains in which input queue {\em addr} is
            written or output queue {\em data} is read are not the external memory
            clock domain then these queues will be asynchronous, i.e. their read
            and write clock domains will be different.
            
            For {\bf read()} access arbitration among competing read and
            write requests is based on a fair strategy (but see the final
            paragraph below). Access requests are funnelled to a single queue
            via a binary tree of queues. At each node conflicting requests
            are resolved in favour of one of the two inputs by a scheme which
            alternates priority between the inputs based on the number of
            tree leaves on each side.
            
            For {\bf readp()} access is subject to access priority, priority
            order being given by input parameter {\bf pri}. Lower numbers are
            higher priority. The priorities do not need to form a contiguous
            set but the same priority must not be used on more than one call
            to {\bf readp()} or {\bf writep()}.
            
            Calls to {\bf read()} and {\bf write()} cannot be mixed with calls
            to {\bf readp()} and {\bf writep()}.


        \subsubsection{class module write - write to external ZBT memory\\
                       class module writep - write to external ZBT memory with priority}\label{sec:lczbtramwrite}

            {\bf Library: zbt\_ram.3pl}

            {\bf Prototype:} \\
            \vsp
            \begin{lstlisting}[gobble=12, language=threepl]
               global module
               write (queue(addrout) <- varargs) {}
            \end{lstlisting}
            {\bf Prototype:} \\
            \vsp
            \begin{lstlisting}[gobble=12, language=threepl]
               global module
               writep (queue(addrout) <- int(pri), varargs) {}
            \end{lstlisting}

            {\bf Output parameters:} \\
            {\bf addrout} is optional and is the address of the datum that
            has been written.

            {\bf Input parameters:} \\
            {\bf pri} is the priority of the write request. \\
            {\bf varargs} expects either two arguments, the address and the data,
            or a single argument which is a struct containing both address and
            data.

            {\bf description:} \\
            The two forms allow the address and data to be given in a single
            queue or as two separate queues or expressions.
            
            For the first form the address and data are contained in a
            struct with fields {\bf address} and {\bf data} that is defined in
            the class -            
            \begin{lstlisting}[gobble=12, language=threepl]
                struct(in_struct,
                    address,    addr_t,
                    data,       data_t,
                    write,      "log"
                );
            \end{lstlisting}
            (the {\bf write} field is assigned internally and can be ignored.)

            
            For the second form the address and data are separate arguments
            which may be any type, but a crude pad/truncate cast to the
            address and data primitive sizes will be performed. At least one of
            the input arguments must contain a queue read.
            
            When the input queue or queues is/are available an external memory
            write will be executed, subject to access arbitration.
           
            If the output parameter is passed a queue variable argument the
            write address is written to the queue when the memory write occurs.
            This can be used to verify that the write has occurred and at what
            address. This might be used for example where the external
            memory is being used as a FIFO and the read address must not
            pass the write address. The output argument may be a queue of type
            "null" if write notification is required but the address value is not
            of interest.
            
            The internal code of this module executes in clock domain {\bf clk}
            (but the call to the module returns with the current clock domain
            restored). If the clock domains in which input queues {\em addr} and
            {\em data} are written are not the external memory clock domain then
            these queues will be asynchronous, i.e. their read and write clock
            domains will be different. Any other sinks to these queues must be in
            the external memory clock domain.
            
            For {\bf write()} access arbitration among competing read and
            write requests is based on a fair strategy (but see the final
            paragraph below). Access requests are funnelled to a single queue
            via a binary tree of queues. At each node conflicting requests
            are resolved in favour of one of the two inputs by a scheme which
            alternates priority between the inputs based on the number of
            tree leaves on each side.
            
            For {\bf writep()} access is subject to access priority, priority
            order being given by input parameter {\bf pri}. Lower numbers are
            higher priority. The priorities do not need to form a contiguous set
            but the same priority must not be used on more than one call to {\bf
            readp()} or {\bf writep()}.
            
            Calls to {\bf read()} and {\bf write()} cannot be mixed with calls
            to {\bf readp()} and {\bf writep()}.


